\begin{textsample}{POS Dim 6 – es – Score 11.00 – es/es\_em\_19.txt}  \label{ex:f6_pos_142}
1.
Visão Geral e Princípios Orientadores Um Projeto Integrador no contexto da Resolução CNE/CEB nº 2/2024 é uma forma de organização pedagógica interdisciplinar que articula componentes curriculares, Itinerários \textbf{Formativos} e a Formação Geral Básica ( FGB ), visando desenvolver competências, habilidades e atitudes a partir de \textbf{problemas} reais.
A Resolução destaca os projetos integradores como estratégia para aprofundar, integrar e contextualizar aprendizagens. 1.1 Princípios Obrigatórios - Interdisciplinaridade e articulação curricular : prioriza a integração entre áreas e componentes. - Contextualização territorial e relevância social : interligar escola e território/lugar e trabalhar \textbf{problemas} reais. - Formação integral e protagonismo estudantil : autonomia, pensamento crítico e Projeto de Vida. - Avaliação \textbf{formativa} e diversificada : registro de evidências ao longo do desenvolvimento. - Planejamento coletivo e gestão democrática : envolvimento de professores, gestão e comunidade. 2.
Etapas do Projeto Integrador 1.
Preparação institucional ( pré-projeto ). 2.
Diagnóstico e escolha do tema. 3.
Concepção e planejamento pedagógico. 4.
Recursos necessários. 5.
Definição de metodologia e \textbf{cronograma}. 6.
Desenvolvimento/Implementação ( execução ). 7.
Registro e monitoramento \textbf{formativo}. 8.
Avaliação final e produto.
Etapa 1.
Preparação institucional Objetivo : garantir condições materiais, de gestão e participação comunitária para viabilizar o projeto, conforme determina a Resolução.
Ações : - Reunião de gestão ( direção + coordenação pedagógica ) para aprovar a proposta inicial. - Formação de um núcleo gestor do projeto : coordenador pedagógico, pedagogo, professor coordenador de área, professores de diferentes componentes, representante de estudantes, representante da família/comunidade ( se possível ). - Verificação de recursos : espaços ( laboratórios, biblioteca ), equipamentos, orçamento, parcerias locais.
Documentar necessidades.
Etapa 2.
Diagnóstico e escolha do tema Objetivo : identificar \textbf{problema} ou tema significativo, pertinente ao território e ao currículo e fazer conexão com a realidade local, promovendo a integração com os Itinerários e a Formação Geral Básica ( FGB ).
Ações : - Levantamento de dados : mapear a realidade local ( realizar entrevistas com a comunidade, pesquisar dados municipais e notícias locais, conversar com especialistas ), levantamento de interesses dos estudantes ( por meio de questionário ou roda de conversa ). - Seleção de temas : gerar lista com temas potenciais.
Critérios de escolha : relevância social, possibilidade de integração curricular, recursos disponíveis e tempo. - Deliberação democrática : votação ou consenso em assembleia de turma e o núcleo gestor.
Deve envolver estudantes e professores. - Exemplo de temas : poluição do rio local ; mapear riscos de trânsito em redondezas ; \textbf{feira} de economia solidária e consumo responsável ; campanha de saúde mental na escola ; museu virtual da memória local.
Etapa 3.
Concepção e planejamento pedagógico Objetivo : produzir o documento do projeto.
Elaborar plano detalhado com objetivos de aprendizagem, competências, áreas envolvidas, produtos esperados, métodos, avaliação e \textbf{cronograma}.
Estrutura sugerida do Documento do Projeto : ver template ( Anexo I ).
Etapa 4.
Recursos necessários - Materiais : ( Liste todos os recursos físicos e \textbf{digitais} ) - Humanos : ( Professores, estudantes, parceiros ) - Financeiros : ( Custos com materiais, transporte, eventos ) Etapa 5.
Definição de metodologia e \textbf{cronograma} operacional Metodologias sugeridas : investigação científica, intervenção social, projetos de extensão, oficinas, seminários, aprendizagem cooperativa, uso de mídias \textbf{digitais}. \textbf{Cronograma} — modelo simples ( 3 meses ) : - Semana 1 : diagnóstico, formalização, levantamento bibliográfico e pergunta-problema. - Semanas 2–4 ( 3 semanas ) : pesquisa aprofundada e ações práticas. - Semanas 5–6 ( 2 semanas ) : sistematização de dados e protótipos ou relatórios. - Semana 7 : revisão, ensaio e preparação da exposição ou publicação. - Semana 8 : evento final ( feira/seminário/exposição ) e avaliação somativa. - Semana 9 : reflexão e registro final ( \textbf{portfólios} ).
Organização de papéis ( exemplo ) : - Coordenador do projeto ( gestor/coordenação pedagógica ) — supervisão geral. - Professor(es) orientador(es) — \textbf{mediação} pedagógica por área. - Estudantes — pesquisadores, comunicadores, técnicos, produtores. - Parceiros externos — apoio técnico e validação ( se houver ).
Ferramentas e recursos \textbf{digitais} : ambientes virtuais ( \textbf{portfólios} \textbf{digitais} ), Google Drive/Office 365, plataformas de apresentação, mídias sociais ( com critérios éticos ), planilhas de monitoramento.
Etapa 6.
Desenvolvimento Princípio : articular teoria e prática ; privilegiar participação ativa dos(as) estudantes ; professor atua como mediador.
Atividades típicas detalhadas : - Oficinas temáticas ( 1–2h ) : trazer especialistas/recursos ; registrar em relatórios. - Investigação de campo : entrevistas, observação, coleta de dados ( fotografias, medidas, questionários ).
Ex. : medir qualidade da água com kits simples. - Aulas integradas : blocos onde conteúdos de diferentes disciplinas trabalham o mesmo subtema ( ex. : cálculo de índices em Matemática ; redação em Linguagens ; leis \textbf{ambientais} em Ciências Humanas ). - Produção de artefatos : boletins, vídeos, protótipos, programas, \textbf{aplicativos}, campanhas, blogs. - Laboratórios de prototipagem / maker : quando projeto envolve \textbf{tecnologia}. - Atividades de \textbf{mediação} socioambiental : intervenção local ( mutirão, campanha educativa ).
Registro obrigatório : \textbf{portfólio} individual e coletivo com evidências ( fotografias, relatórios, diários de campo, produções ).
A Resolução enfatiza documentação das evidências para avaliação.
Etapa 7.
Monitoramento \textbf{formativo} ( acompanhamento ) Como fazer : - Reuniões quinzenais do núcleo gestor para acompanhar andamento. - Relatórios de progresso por grupo ( modelo : o que foi feito, dificuldades, próximos passos ). - Registros de aprendizagem : fichas de observação, autoavaliação dos estudantes, tutoriais. - Intervenção pedagógica quando há estudantes em risco ( apoios, reforço ).
Instrumentos : planilha de acompanhamento, diário de campo dos professores, formulário de feedback dos estudantes.
Etapa 8 — Avaliação final e produto Princípios de avaliação : avaliação deve ser \textbf{formativa} e somativa, diversificada e registrada.
Componentes da avaliação final ( sugestão ) : - Produto ( 40 \% ) : qualidade técnica, pertinência, inovação, impacto. - Processo/portfólio ( 30 \% ) : registro de atividades, evidências de pesquisa e desenvolvimento, participação. - Apresentação pública ( 15 \% ) : clareza, argumentação, uso de \textbf{mediações}. - Avaliação entre pares e autoavaliação ( 10 \% ) : reflexão sobre contribuição individual e coletiva. - Relatório docente/observações ( 5 \% ) : avaliação \textbf{qualitativa} do professor orientador.
Modelo de rubrica ( resumida ) — Produto ( 40 \% ) - Relevância do \textbf{problema} : 0–10 - Fundamentação teórica/metodológica : 0–10 - Qualidade técnica ( projeto/produto ) : 0–10 - Impacto/viabilidade de intervenção : 0–10

% matched lemmas: ambiental, aplicativo, cronograma, digital, feira, formativo, mediação, portfólio, problema, qualitativo, tecnologia
\end{textsample}
